\subsection{Loss releases and finance retain different origins}
\label{sec:insurance-service-finance-origins}

The single-service-period example releases the whole loss component at once.
Across several periods, finance also changes the liability before services
release it. Paragraphs 50--52 require a systematic allocation of claim releases,
risk-adjustment releases and insurance finance between the loss component and
the remainder of the liability. The loss component must reach zero by the end
of coverage \citep[IFRS~17.49--52]{eu20231803}. Allocating only service releases
can leave an unexplained balance when finance has changed the amount to release.

Consider a separate group with remaining-coverage liability 100, loss component
20 within that liability, and no CSM. Cash is 200 and equity is 100. Assume
no new estimates, investment components, acquisition-cost presentation or tax,
and present all finance in profit. For this group, supply a supported fixed
proportional allocation of one fifth. Let \(L\) denote the liability and
\(K\) its loss component, so \(k=K/L=1/5\). Positive finance expense
\(h\) increases the liability; expected claims and risk-adjustment releases
\(c\) and \(a\) reduce it. Under this declared allocation,
\[
 L'=L+h-c-a,\qquad
 K'=K+k h-kc-ka=kL'.
\]
The last equality explains why exhausting the supported liability releases also
exhausts the loss component. It does not prescribe proportional allocation for
every group. The finance allocation \(kh\) remains a finance movement; the
service reversal is \(k(c+a)\). Revenue from these releases is
\((1-k)(c+a)\), excluding that loss reversal.

In the first period, finance expense 10 raises liability to 110 and its loss
component to 22. Expected claim release 50 and risk-adjustment release 5 then
release 11 from the loss component and produce revenue 44. Actual claims of
50 are incurred and paid separately. Liability closes at 55, loss component
at 11, cash at 150 and equity at 95. Service profit is
\(44+11-50=5\); finance expense 10 makes period profit \(-5\).

In the second period, finance income 5 reduces liability to 50 and the loss
component to 10. Claim release 45 and risk-adjustment release 5 exhaust both:
revenue is 40 and loss reversal is 10. Paying actual claims 45 leaves cash
and equity of 105. Service profit is again 5; finance income makes this period's
profit 10. Across both periods, service profit 10 less net finance expense 5
gives total profit 5. Treating all releases of 105 as revenue and then adding
loss reversals of 21 would overstate profit by 21. The separate categories,
rather than final cash alone, reveal the mistake.

The held treaty needs the same discipline. Recall the covered-loss example:
issued liability 120 contains loss component 20, while a timely half-quota
treaty covers contracts with loss 8. Held fulfilment asset 24.5 and cost CSM
6 give asset 30.5, containing a loss-recovery component of 4. With cash 123.5,
equity is 34. Reducing covered expected claims by 4 reduces issued liability
to 116 and loss component to 16. Held recovery falls by 2, giving held asset
28.5, unchanged cost CSM 6 and loss-recovery component 2. Equity rises to 36.
Paragraph 66(c)(i) puts the held fulfilment change in profit because the related
issued change did not adjust issued CSM. Paragraph 66(bb) addresses reversals
\emph{not already included} in held fulfilment changes; the same 2 cannot
be charged again \citep[IFRS~17.66(bb),66(c)(i),B119F]{eu20231803}.
Subsequent issued and held service, claims of 114 and recoveries of 22 close
at cash 31.5, as before. Tracking the recovery component does not create a
second asset or a second expense.

\subsection{A change in holdings leaves an old OCI schedule}
\label{sec:insurance-frozen-oci-origins}

Current insurance finance and the reclassification of old OCI answer different
questions. Under paragraph 89(b), profit finance on a direct-participation group
matches the profit income or expense of the underlying items actually held.
Suppose total insurance finance expense is \(H\), and those items produce
profit income \(A\). Insurance profit expense is then \(A\), while
insurance OCI income is \(A-H\). Asset income and insurance expense remain
separately presented, with net profit finance zero
\citep[IFRS~17.89--90,B134]{eu20231803}.

When a holding change moves the insurer between 88(b) and 89(b), B135--136
requires previously accumulated OCI to be reclassified using the former
approach and the assumptions immediately before the change. Those assumptions
are frozen, and comparatives are unchanged \citep[IFRS~17.B135--136]{eu20231803}.
The new policy governs new-period finance. Recalculating the old balance as if
the new policy had always applied would erase the required history.

Use the terminal-payment arithmetic from the preceding variable-benefit example
in a separate, supplied direct-participation case with no remaining service,
CSM or nonfinancial risk adjustment. At the holding change, payment 144 is due
in two years. Its frozen current value at 10\% is \(14400/121\), while
the old systematic carrying amount is 100 with allocation rate 20\%.
Old insurance OCI is therefore
\(100-14400/121=-2300/121\). Qualification and actual holdings are supplied
findings; these values cannot establish either on their own.

Under the old assumptions, after one year the current-value calculation would
be \(1440/11\) and the systematic amount would be 120. Old OCI would be
\(-120/11\), reaching zero at payment. Thus the first year's increase in
old OCI, and equal profit expense from reclassification, is
\[
 -\frac{120}{11}-\left(-\frac{2300}{121}\right)
 =\frac{980}{121}.
\]
The second-year reclassification expense is \(120/11\). These amounts sum
to the original expense \(2300/121\). They do not depend on the next
period's actual current valuation.

To see both histories together, suppose the insurer now holds a matching
underlying asset initially measured at \(14400/121\), with post-acquisition
cash of 200. The acquisition cash is part of this supplied opening snapshot.
In the first year, actual asset and liability both become 130. The asset
reports profit income \(1440/121\) and OCI expense \(10/11\), which
together reconcile its increase. New insurance profit finance exactly offsets
that asset profit income; new insurance OCI is income \(10/11\), offsetting
the asset OCI. The old \(980/121\) reclassification remains a separate
profit expense and equal OCI income. Total equity is unchanged, although
profit and OCI each contain that old history.

If the actual first-year asset and liability instead become 132, with the same
asset profit income, asset OCI becomes income \(12/11\) and new insurance
OCI becomes expense \(12/11\). The old reclassification is still
\(980/121\). A method that changes this old amount in response to the new
valuation has failed the frozen-assumption requirement. At terminal receipt
and payment of 144, the asset, liability and both OCI schedules are exhausted;
cash and equity remain 200.

The reverse direction also needs a former-policy calculation. Suppose a held
underlying FVOCI asset has fair value 121, amortised cost 100 and original
effective rate 10\%, with a terminal payment 121 in two years. Under the
supplied frozen zero-current-rate assumptions, its fair value remains 121
while amortised cost grows to 110 and then 121. A matching insurance history
therefore has old OCI \(-21,-11,0\). After the holdings change to 88(b),
the old balance reclassifies as expenses 10 and 11. The new systematic
allocation begins at 121 with rate zero. If the actual liability becomes 130
and then returns to 121, new OCI is expense 9 and then income 9; it does not
replace the old 10 and 11 schedule. Post-change cash 200 and liability 121
give equity 79, which returns to cash and equity 79 after settlement. The
actual holdings change and asset accounting remain independently supplied.

\subsection{Derecognition follows the origin of accumulated OCI}
\label{sec:insurance-origin-derecognition}

A transfer under paragraph 91 can have identical cash and total equity but
different profit presentation. Take cash 200, liability 100 and equity 100,
including accumulated insurance OCI expense 12. Assume a qualifying transfer
at carrying, with no remaining service, CSM, risk adjustment or pending
policy-change runoff. Paying 100 and removing the liability leaves cash and
equity 100.

If that OCI arose under 88(b), paragraph 91(a) reclassifies it to profit:
expense 12 in profit and income 12 in OCI, leaving no accumulated balance.
If it arose under 89(b), paragraph 91(b) prohibits that reclassification
through profit \citep[IFRS~17.91]{eu20231803}. In this finite presentation,
the expense remains within equity in OCI; no separate within-equity transfer
is assumed. A positive accumulated OCI amount reverses these signs but not
the policy distinction. Selecting a rule from the current policy label alone
would discard the fact that paragraph 91 needs.

The records consequently retain each balance's original policy, actual holding
history and any frozen reclassification assumptions. They distinguish finance
measurement, current presentation and old-balance reclassification before
forming profit and total equity. Mixed-origin early termination during an
unfinished policy-change schedule still requires a specific source-supported
disposition; these homogeneous transfers do not settle that question.
